\documentclass{article}
\usepackage[utf8]{inputenc}
\usepackage[T1]{fontenc}
\usepackage{geometry}
\geometry{
    a4paper,
    left=2.5cm,
    right=2.5cm,
    top=2.5cm,
    bottom=2.5cm
}
\usepackage{graphicx}
\usepackage{float}
\usepackage{hyperref}
\usepackage[english]{babel}
\usepackage{amsmath}
\usepackage{amssymb}
\usepackage{array}
\usepackage{booktabs}
\usepackage{caption}
\usepackage{subcaption}
\usepackage{enumitem}
\usepackage{fancyhdr}
\usepackage{titlesec}
\usepackage{url}
\usepackage{textcomp}
\usepackage{xcolor}
\usepackage{listings}
\usepackage{multirow}
\usepackage{longtable}
\usepackage{tabularx}
\usepackage{times}

% Set up page style
\pagestyle{fancy}
\fancyhf{}
\rfoot{\thepage}
\renewcommand{\headrulewidth}{0pt}

% Title formatting
\titleformat{\chapter}[display]{\normalfont\bfseries}{\chaptertitle\thechapter}{0pt}{\Huge}
\titlespacing*{\chapter}{0pt}{50pt}{40pt}

% Document info
\title{Template Document}
\author{Academic Support Agent}
\date{\today}

\begin{document}
\maketitle

\maketitle
\thispagestyle{fancy}
%\section{Contexte, positionnement et objectif(s) de la proposition}
\section{Context, position and objectives of the proposal}
%Sur une page et demie (1,5) environ, 
%D crire les objectifs et les hypoth ses scientifiques
%Montrer l originalit  et la pertinence par rapport   l tat de l art
%D crire la m thodologie et la gestion des risques

%\section{Organisation du projet et moyens mis en  uvre}
\section{Project organisation and implementation}

%Sur une demi (0,5) page environ,
%S il s agit d un projet collaboratif (PRC, PRCE), 
%Pr senter succinctement le coordinateur scientifique et son implication :
%Indiquer l exp rience du coordinateur scientifique et son r le dans le projet
%Pr senter le consortium et sa compl mentarit  : 
%Indiquer les diff rentes comp tences envisag es pour mener le projet objet de la pr -proposition, en pr cisant l identit  du ou des scientifiques impliqu s, l identification des  tablissements auxquels ils sont rattach s et 2-3 r f rences pertinentes dans le domaine en lien direct avec la pr -proposition (publications, faits marquants de R&D, brevets, prix scientifiques, produits, proc d s, licences, services...), et tout autre  l ment permettant de juger de la qualit  des partenaires et du consortium le cas  ch ant. 
%D crire succinctement les moyens mis en  uvre pour atteindre les objectifs :
%Indiquer notamment le montant total de l aide pr visionnelle demand e, sa r partition entre partenaires et entre grands postes de d penses ( quipement, personnel, prestations, autre fonctionnement) et les moyens de recherche compl mentaires au projet.
%
%S il s agit d un projet jeune chercheuse -- jeune chercheur (JCJC), 
%Pr senter succinctement le coordinateur scientifique, sa situation au moment de la soumission et lors de la mise en  uvre du projet, son r le dans le projet 
%Pr senter l quipe du coordinateur scientifique et son positionnement scientifique au sein de l organisme ou du laboratoire d accueil.
%D crire succinctement les moyens mis en  uvre pour atteindre les objectifs :
%Indiquer notamment le montant total de l aide pr visionnelle demand e, sa r partition entre grands postes de d penses ( quipement, personnel, prestations, autre fonctionnement) et les moyens de recherche compl mentaires au projet.

%\section{Impact et retomb es du projet}
\section{Project impact and benefits}
%Sur une (1) page environ,
%D crire dans quel(s) domaine(s) (scientifique,  conomique, social ou culturel) les r sultats du projet peuvent avoir un impact
%D crire en quoi les attendus du projet r pondent aux enjeux d un (axe de) d fi du Plan d Action 2017 de l ANR, de la SNR, de la recherche fondamentale
%D crire la strat gie de diffusion et/ou de valorisation envisag e : l objectif est d indiquer comment les r sultats seront diffus s et/ou valoris s (publications scientifiques ou grand public, brevets, etc.), et non de d tailler les modalit s de partage de l ventuelle propri t  intellectuelle entre partenaires

%%%% Couper le PDF ici %%%%
\pagebreak
\section*{\bibname}
\printbibliography[heading=none]
\end{document}